\documentclass[twocolumn]{aastex701}
\begin{document}


\title{Variable stars hidden from Gaia DR3: a ZTF census with a pre-registered prediction for Gaia DR4}

\author[0000-0001-8779-0700]{Deepnil Ray}
\affiliation{Ashoka University, Sonipat, Haryana 131029, India}
\email{deepnilray_ug2025@ashoka.edu.in}
\affiliation{The NoRCEL Institute, United Kingdom}

\begin{abstract}
Gaia DR3 published variability results for about 10.5 million sources and none for the rest. I searched ZTF DR23
light curves of 13.9 million stars that DR3 left without any variability result, using a detector frozen and
validated on sky tiles sealed before any data were opened (AUC 0.954; false-positive rate 0.23\%). The search yields
5,778 periodic variables confirmed in independent light curves and stable season by season; the ZTF SCoPe
classifications already list 86.8\% of them as variable, an independent confirmation that the detections are real.
Thirty-two stars are listed as variable in no published catalogue and pass sky-image and light-curve checks; they
are new variables, including three detached eclipsing binaries (orbital periods 1.89, 3.02 and 25.4 d). Before Gaia
DR4 (2 December 2026) I fix and publish a prediction of what DR4 will report for these stars against matched
controls, turning DR4 into an independent test of the method and a measurement of which variables Gaia misses.
\end{abstract}

\keywords{Variable stars (1761) --- Sky surveys (1464) --- Astronomy data analysis (1858)}

\section{Motivation}
Gaia's sparse sampling, set by its scanning law, leaves real variables unflagged, but how many, and of which
periods and amplitudes, has not been measured; \citet{Everall2022} describe a selection function for Gaia's
variable samples as ``hugely valuable'' and note that building one requires knowing which sources Gaia would have
classified as variable had they been variable. A census of real variables among stars Gaia called quiet answers
that question directly. Testing such a census against the next release, with the prediction fixed in advance,
makes the answer verifiable.

\section{Data and method}
I selected Gaia DR3 stars with no \texttt{vari\_summary} row, not flagged as variable, and RUWE $<1.4$, at
$|b|\ge15^\circ$ and $\delta\ge-25^\circ$, and read their ZTF DR23 light curves from the IRSA HATS archive. A
seeded manifest assigned every HEALPix tile to training, validation or search before any download. A gradient-boosted
classifier on light-curve shape features (Lomb--Scargle power and false-alarm probability, amplitude over error,
reduced $\chi^2$, skewness, von Neumann ratio and an eclipse box-fit channel) was frozen before the validation tiles
were opened. On 356 Gaia-classified variables and 4,000 Gaia-quiet stars in those tiles it reached AUC 0.954
(lower 95\% bound 0.939; shuffled-label control 0.494). The full search rule has a false-positive rate of 0.23\%
(upper bound 0.43\%), itself an upper limit because some ``negatives'' are unrecognised variables. Every candidate
had to recover its period in at least two independent ZTF light curves, in at least half of the seasons tested,
away from sampling aliases, with no brighter Gaia star within $6''$. Only grades A and B are used.

\section{Results}
Table~\ref{tab:res} summarises the census. Prior art was checked against VSX, SIMBAD, ALeRCE, 17 VizieR catalogues
and the latest SCoPe release \citep{Healy2024}, which covers 5,760 of the 5,778 stars. SCoPe lists 86.8\% as variable
with a matching class or period, independently confirming that the detections are real; it lists 11.7\% below its
own variability threshold and 1.5\% are absent from it. Ninety stars are listed as variable in no catalogue. Each
was inspected on Pan-STARRS1 images and its light curve by eye: two are blends and were dropped; 32 were held back
because of a possible blend, a known variable within $30''$ whose light may contaminate the ZTF photometry, or ZTF
magnitudes inconsistent with Gaia; 24 have light curves too weak to claim. The remaining 32 are new variable stars
(Table~\ref{tab:new}). They include three detached eclipsing binaries, one with a 25.4-d orbit, and two contact
binaries near 0.25--0.29 d. The ZTF DR1 variable candidates of \citet{Ofek2020} could not be checked and are a stated limitation.

\begin{deluxetable}{lc}
\tablecaption{Census of Gaia-DR3-quiet variables at the catalogue freeze (27 November 2026, 03:45 UTC).\label{tab:res}}
\tablehead{\colhead{Quantity} & \colhead{Value}}
\startdata
Sky tiles searched ($|b|\ge15^\circ$) & 2,796 (9,400 deg$^2$) \\
Gaia-DR3-quiet stars with ZTF light curves & 13.9 million \\
Validated A/B periodic variables & 5,778 \\
\quad in SCoPe as variable, matching class or period & 5,000 (86.8\%) \\
\quad in SCoPe below its variable threshold & 676 (11.7\%) \\
\quad absent from SCoPe & 84 (1.5\%) \\
Listed as variable in no catalogue & 90 \\
\quad new variables (clear light curve) & 32 \\
\enddata
\end{deluxetable}

\section{Prediction for Gaia DR4}
The following is fixed before DR4 appears; the files and their hashes are in the Zenodo record.
(P1) Among stars matched through DR4's \texttt{dr3\_neighbourhood} table, DR4 flags a larger fraction of the 5,778
candidates as VARIABLE than of 550 control stars matched in position (within $0.5^\circ$), $G$ ($\pm0.25$) and
$G_{BP}-G_{RP}$ ($\pm0.15$), at $z>3.3$. (P2) DR4 G-band light curves show the ZTF period (factors 0.5, 1 or 2, within 1\%)
with false-alarm probability $<10^{-3}$ more often for candidates than for controls tested at their paired period.
(P3) Among candidates, the DR4 VARIABLE fraction rises with amplitude and falls for faint stars and periods beyond
$\sim$10 d, measured in bins fixed now (8 in $\log P$ over 0.05--500 d, 5 in amplitude, 1-mag bins in $G$ from 12 to 19).
No target fraction is set; the outcome will be reported whatever it is.

\section{Discussion}
A census fixed before a data release and then tested by it is a strong check that the detections are real. The
stars DR4 still fails to flag will define where Gaia's variability detection is blind, the input a variability
selection function needs \citep{Everall2022}. The DR4 result and the full selection function in period, amplitude,
magnitude and scanning-law coverage will follow in a forthcoming paper.

\startlongtable
\begin{deluxetable*}{lrrrrrc}
\tablecaption{The 32 new variable stars. $P$ is the orbital period for EA/EW systems; for the others it is the ZTF
period, and the true period may be twice it. Amp.\ is the ZTF $r$-band amplitude.\label{tab:new}}
\tablehead{\colhead{Gaia DR3 source\_id} & \colhead{RA (deg)} & \colhead{Dec (deg)} & \colhead{$G$} & \colhead{$P$ (d)} & \colhead{Amp.} & \colhead{Type}}
\startdata
220985864201130880 & 52.68684 & +33.99985 & 16.77 & 0.34235 & 0.16 & -- \\
144104918810873600 & 44.90387 & +41.50692 & 14.67 & 0.48126 & 0.09 & -- \\
3212111109018082688 & 77.52603 & -4.16784 & 17.02 & 0.80700 & 0.13 & -- \\
1115128064441569152 & 101.55789 & +74.76635 & 16.31 & 0.80815 & 0.15 & -- \\
4542539793743461376 & 262.44965 & +13.98629 & 16.50 & 0.81948 & 0.11 & -- \\
1198491283713804928 & 241.74909 & +16.51748 & 17.12 & 1.89020 & 0.16 & EA \\
1052374297278903936 & 158.41634 & +63.36611 & 16.56 & 3.02409 & 0.13 & EA \\
568272281191898240 & 57.94509 & +81.48817 & 14.97 & 6.26475 & 0.13 & -- \\
321766687685841536 & 19.20328 & +36.97374 & 15.27 & 22.02554 & 0.14 & -- \\
346276760375634816 & 31.08684 & +42.72787 & 16.32 & 25.44693 & 0.08 & EA \\
1737667592423979520 & 313.02030 & +7.58462 & 18.10 & 0.24313 & 0.19 & -- \\
4230919197495489024 & 306.25634 & +1.40315 & 18.10 & 0.24756 & 0.16 & EW \\
883343449870017152 & 107.85309 & +26.40558 & 17.11 & 0.28981 & 0.12 & EW \\
905894192835199744 & 122.34616 & +35.11849 & 18.32 & 0.39661 & 0.32 & -- \\
2091030666583675776 & 279.85317 & +32.80090 & 18.07 & 0.48658 & 0.18 & -- \\
392034654752728064 & 5.22944 & +46.50365 & 17.52 & 0.56810 & 0.16 & -- \\
1788985957620546432 & 318.80600 & +19.85998 & 17.48 & 0.58055 & 0.16 & -- \\
2140005873718290688 & 290.79697 & +54.26231 & 15.42 & 0.78431 & 0.07 & -- \\
4475953030265090560 & 268.22324 & +7.72803 & 17.14 & 0.84516 & 0.16 & -- \\
370109740058043776 & 18.90870 & +38.88437 & 17.12 & 0.91009 & 0.13 & -- \\
4586091900291068672 & 273.41947 & +26.89792 & 16.85 & 0.97166 & 0.09 & -- \\
385049354301256704 & 0.29195 & +44.66341 & 17.00 & 1.20782 & 0.13 & -- \\
4536646724957687168 & 276.53158 & +24.82873 & 16.64 & 1.28410 & 0.08 & -- \\
2269181859866703232 & 272.11790 & +75.81729 & 17.06 & 2.53832 & 0.12 & -- \\
309225245742174464 & 21.69344 & +29.75539 & 16.70 & 2.92022 & 0.10 & -- \\
3229653954478382208 & 71.87815 & -0.96564 & 13.33 & 4.33721 & 0.06 & -- \\
2298065927249128704 & 306.26974 & +80.75639 & 17.38 & 4.84942 & 0.12 & -- \\
165986334076500096 & 63.54724 & +30.73295 & 17.05 & 5.44164 & 0.15 & -- \\
3204803754739717888 & 67.16010 & -3.26040 & 14.29 & 6.41849 & 0.11 & -- \\
1273135410136152960 & 233.89113 & +29.72690 & 13.76 & 6.72214 & 0.12 & -- \\
4584902503592643200 & 276.03126 & +26.00164 & 18.45 & 10.39817 & 0.17 & -- \\
340135777480090112 & 39.24818 & +42.95679 & 17.03 & 10.90664 & 0.10 & -- \\
\enddata
\end{deluxetable*}

\section*{Data availability}
The frozen catalogue, the pre-registered prediction and the code are archived on Zenodo (DOI to be added; deposited 27 November 2026) and at \url{https://github.com/deepnilray/SearchChev}.

\begin{acknowledgments}
This work used ZTF data (IRSA) and data from the ESA mission Gaia, processed by the Gaia DPAC. Software development
and analysis were assisted by an AI system (Claude, Anthropic); all results were checked by the author.
\end{acknowledgments}

\begin{thebibliography}{}
\bibitem[Everall \& Boubert(2022)]{Everall2022} Everall, A., \& Boubert, D. 2022, MNRAS, 509, 6205
\bibitem[Healy et al.(2024)]{Healy2024} Healy, B. F., et al. 2024, ApJS, 272, 14
\bibitem[Ofek et al.(2020)]{Ofek2020} Ofek, E. O., et al. 2020, MNRAS, 499, 5782
\end{thebibliography}

\end{document}
